\section{Conclusion and Formal Limitations}
\label{sec:conclusion}

In this paper, we have presented a formal configuration research model for \textbf{Sod Ha-Bechinah} (סוד הבחינה).

We have established:
\begin{enumerate}
    \item The core relational model $B = \operatorname{Bechina}(V, R, S)$.
    \item The four fundamental dimensions $\mathcal{D} = (P, C, R, \tau)$ and information encoding layers $\mathcal{I} = (T, N, G, L)$.
    \item Hierarchical state vectors $\Omega_i = (V_i, R_i, S_i, T_i, N_i, G_i, L_i)$ with refinement maps $\Phi_i$ and projection maps $\Pi_i$.
    \item Mathematical formalizations for Adam Kadmon Gematria sets $\mathcal{H} = \{72, 63, 45, 52\}$, Akudim/Nekudim/Berudim transformations, Sefirotic graph networks, and the 231 relational gates $G_{231} = \binom{22}{2}$.
\end{enumerate}

\subsection*{Formal Limitations}
We conclude with an explicit methodological boundary:
\begin{quote}
This project does not claim to mathematically prove metaphysical or theological propositions. Instead, it constructs a formal, self-consistent modeling language in which relations, scales, configurations, information layers, transformations, and hierarchical visibility can be represented rigorously and implemented computationally.
\end{quote}
